\chapter*{Conclusion}
\addcontentsline{toc}{chapter}{Conclusion}

La régression logistique repose sur une chaîne cohérente : modèle linéaire des
log-cotes, loi de Bernoulli conditionnelle, maximum de vraisemblance, log-perte
convexe et optimisation numérique contrôlée. OvR réutilise ce problème binaire
$K$ fois et fournit une règle multiclasses par maximum des scores. Sa simplicité
a une contrepartie précise : la somme des sorties indépendantes prend une valeur
libre, quand une distribution catégorielle vaut $1$. La validité du résultat
tient enfin autant au protocole expérimental --- séparation stricte du test,
diagnostics par classe, archivage --- qu'à la correction des équations.

\begin{resultat}[Enchaînement des opérations]
Le modèle reçoit des caractéristiques préparées, calcule quatre combinaisons
linéaires et retient le score maximal. Son apprentissage choisit chaque vecteur
de paramètres de façon à maximiser la vraisemblance de la cible binaire associée.
Sa validation porte sur trois questions séparées : la convergence numérique de
l'optimisation, la généralisation hors échantillon, et les erreurs propres à
chaque maison. Chacune se mesure par ses propres quantités.
\end{resultat}
